\section{Introduction}

One of the early consequences of the Axiom of Choice $\AC$ was the
existence of bases for vector spaces. In 1905, shortly after Zermelo's
proof of the well-ordering theorem~\cite{Zermelo1904}, Hamel constructed
a basis of $\mathbb R$ over $\mathbb Q$ and used it to obtain
discontinuous solutions of the functional equation
$f(x+y)=f(x)+f(y)$~\cite{Hamel1905}.
Hausdorff's treatment of abstract real linear spaces in
1932 included the existence of bases~\cite[p.~295]{Hausdorff1932}.
For vector spaces over arbitrary fields, one of the earliest
recorded proofs, due to Zorn, appears in 1935~\cite{Zorn1935} in his
treatment of what is modernly known as Zorn's lemma.

Thus, $\AC$ implies that every vector space has a basis. Whether the
existence of bases really depended on the Axiom of Choice or whether it
could be proved in the Zermelo--Fraenkel set theory $\ZF$ became a more
complicated question.

In 1962, L\"auchli~\cite{Lauchli1962} constructed permutation models
with self-singletons $x=\{x\}$ containing a vector space with no basis,
and other such models containing a vector space with bases of different
cardinalities. Since both Choice and Foundation fail in these models,
his constructions established the independence of the existence of
bases from $\ZFm$, but did not determine the precise choice principle
implied by the assertion that every vector space has a basis.

Some progress was first obtained for stronger assertions about bases.
Working in $\ZF$, Bleicher~\cite{Bleicher1964} derived the axiom of
multiple choice $\MC$ from the assertion that every linearly independent
set extends to a basis. Since $\MC$ is equivalent to $\AC$ in
$\ZF$~\cite{FelgnerJech1973}, this extension principle implies $\AC$.
The equivalence between $\MC$ and $\AC$ uses Foundation: L\'evy's
permutation-model construction establishes the relative consistency of
$\MC+\neg\AC$ with $\ZFm$~\cite{Levy1962}. Halpern~\cite{Halpern1966}
later proved directly in $\ZFm$ that $\AC$ follows from another stronger
assertion, namely that every generating set of every vector space
contains a basis. Neither result, however, settled the question for the
weaker assertion that every vector space has some basis.

Blass made the decisive advance in 1984~\cite{Blass}. Working in
$\ZFm$, he showed that the existence of bases for all vector spaces
implies $\MC$.

\begin{theorem}[Blass {\cite{Blass}}]\label{thm:blass}
In $\ZFm$, if every vector space over a field of characteristic zero has a basis, then $\MC$ holds.
Thus,
\[
    \ZF\vdash \text{``every vector space has a basis''} \leftrightarrow \AC.
\]
\end{theorem}

The restriction to fields of characteristic zero is not stated in
Blass's article, but it follows from the proof given there.

Blass's theorem settles the converse in $\ZF$ because, as observed
above, $\MC$ implies $\AC$ in $\ZF$. Without Foundation, it yields only
$\MC$. The remaining question was therefore whether, in $\ZFm$, the
existence of bases implies $\AC$
itself~\cites[p.~33]{Blass}{Karagila2017MO}[p.~423]{Morillon}[p.~121]{Philip2025}.

The apparent role of Foundation in this problem is particularly striking,
since Foundation is usually regarded as having little bearing on ordinary
mathematics.
Jech writes that this axiom is ``irrelevant for the development of ordinal
and cardinal numbers, natural and real numbers, and in fact of all ordinary
mathematics'' while emphasizing its usefulness in the metamathematics
of set theory~\cite[Chapter~6, p.~63]{Jech2003}.
Kunen writes: ``The Axiom of Foundation is, as always in mathematics,
totally irrelevant''~\cite[p.~47]{Kunen1980}.
Similar remarks appear in
\cites[Chapter~III, \S4, p.~101]{Kunen1980}[\S I.14]{Kunen2009}[p.~87]{FraenkelBarHillelLevy1973}.
Nevertheless, Foundation is essential to the implication
$\MC\Rightarrow\AC$ used in Blass's proof.
Determining whether Foundation can be eliminated from the equivalence
between the existence of bases and $\AC$ therefore clarifies whether this axiom
plays an essential role in the proof or whether the existence
of bases alone suffices.

In this paper we settle this question by showing that the existence
of bases is enough to derive $\AC$.
\begin{theorem}\label{thm:main}
In $\ZFm$, assume that every vector space over a field of characteristic
zero has a basis. Then $\AC$ holds.
Thus, in $\ZFm$, the following are equivalent:
\begin{enumerate}
    \item $\AC$;
    \item every vector space has a basis;
    \item every vector space over a field of characteristic zero has a basis.
\end{enumerate}
\end{theorem}

This paper is organized as follows.

In Section~\ref{sec:preliminaries}, we fix set-theoretic notation and recall some
elementary set-theoretic facts.

In Section~\ref{sec:galois}, we review some basic field extension theory
that will be used in the proof of Theorem~\ref{thm:main}.
All the statements in this section are standard and are reviewed
for the reader's convenience, to fix notation and to ensure that
the proofs do not require Foundation or Choice, although the proofs are standard.

In Section~\ref{sec:finite}, we prove Lemma~\ref{lem:finite-bases}, a criterion for the finiteness of a normal separable algebraic extension in terms of polynomial coefficients associated
to finite collections of bases and intermediate finite-dimensional extensions.

In Section~\ref{sec:extension}, we use Lemma~\ref{lem:finite-bases} to prove Theorem~\ref{thm:algebraic}, which shows that, assuming $\MC$ and given a field extension $L/K$ that is normal, separable and algebraic, if certain intermediate extensions of $L/K$ have bases and if $K$ is well-orderable, then this well-ordering can be extended to a well-ordering of $L$.

In Section~\ref{sec:main-proof}, we prove Theorem~\ref{thm:main} by showing that every set can be embedded into an algebraic, normal and separable extension of a well-orderable field of characteristic zero and then applying Theorem~\ref{thm:algebraic} to well-order this extension.

Finally, in Section~\ref{sec:conclusion}, we discuss some consequences of Theorem~\ref{thm:main} and make final remarks.
